% Generated draft: requires mathematical and source review.
\documentclass[11pt]{article}
\usepackage[margin=25mm]{geometry}
\usepackage{fontspec}
\setmainfont{Latin Modern Roman}
\usepackage{amsmath,amssymb,mathtools,unicode-math}
\setmathfont{Latin Modern Math}
\usepackage{enumitem,needspace,xurl,hyperref}
\setlength{\parindent}{0pt}
\setlength{\parskip}{5pt}
\setlength{\emergencystretch}{3em}
\newcommand{\N}{\mathbb{N}}
\newcommand{\Z}{\mathbb{Z}}
\newcommand{\Q}{\mathbb{Q}}
\newcommand{\R}{\mathbb{R}}
\newcommand{\C}{\mathbb{C}}
\newcommand{\F}{\mathbb{F}}
\newcommand{\A}{\mathbb{A}}
\newcommand{\PP}{\mathbb{P}}
\newcommand{\E}{\mathbb{E}}
\newcommand{\eps}{\varepsilon}
\newcommand{\dd}{\,\mathrm{d}}
\begin{document}
\section*{UnsolvedMath ID: 30006094}
\textbf{Source problem code:} catalogue entry
\subsection{Definitions and mathematical statement}
The supplied record is an expository fragment, not a recoverable standalone problem. No definite conjecture or question can be stated from it.

Let \(\mathrm{Sp}\) denote the category of spectra, and let \(\Sigma\) denote suspension. For a spectrum \(E\), the fragment denotes Bousfield localization with respect to \(E\) by
\[
L_E\colon \mathrm{Sp}\longrightarrow\mathrm{Sp},
\]
and its image category by \(L_E\mathrm{Sp}\). It specializes to a spectrum denoted \(K(n)\); write \(K(n)_*(X)\) for the associated graded homology of a spectrum \(X\). The excerpt does not define \(K(n)\), the spectra \(V(n)\), or the elements \(v_n\).

With \(d_n=|v_n|\) denoting the unspecified suspension degree, the corrupted display appears to describe a cofiber sequence
\[
\Sigma^{d_n}V(n-1)\xrightarrow{w_n}V(n-1)\longrightarrow V(n).
\]
This is a tentative reconstruction of the extracted notation. The fragment gives the parameter restrictions
\[
0\leq n\leq 3,\qquad p\geq 2n+1,
\]
but the placement of this qualification is disrupted. Its precise scope, the domain of \(p\), and the meaning of \(V(-1)\) at \(n=0\) are not specified. In particular, a prime domain for \(p\) must not be silently imposed.

The fragment asserts that the induced map
\[
K(n)_*(w_n)\colon K(n)_*\bigl(\Sigma^{d_n}V(n-1)\bigr)
\longrightarrow K(n)_*\bigl(V(n-1)\bigr)
\]
is an isomorphism, that all iterates of \(w_n\) are essential, and that the telescope
\[
\operatorname*{hocolim}\left(
V(n-1)
\xrightarrow{\Sigma^{-d_n}w_n}\Sigma^{-d_n}V(n-1)
\xrightarrow{\Sigma^{-2d_n}w_n}\Sigma^{-2d_n}V(n-1)
\longrightarrow\cdots
\right)
\]
is noncontractible. Here essential means not nullhomotopic, and the telescope maps are interpreted as the indicated desuspensions of \(w_n\). These are supplied expository assertions, not an independently verified theorem statement. The excerpt does not explicitly supply the nonvanishing hypothesis on \(K(n)_*(V(n-1))\) needed to justify its stated inference from an isomorphism to essentiality and noncontractibility.

The excerpt also qualitatively contrasts computing \(\pi_*L_{K(2)}V(1)\) with computing \(\pi_*V(1)\); it supplies neither a computation nor a precise computability claim. Despite the heading concerning the telescope conjecture, no formulation of that conjecture or requested mathematical conclusion is present.
\subsection{Short English statement}
This record contains damaged background exposition about localization, a proposed cofiber sequence, and its telescope. It does not specify a standalone problem. Its parameter conventions and some hypotheses remain unresolved, so the displayed reconstruction requires human review.
\subsection{Sources}
\begin{itemize}
\item Local UnsolvedMath registry: supplied complete record, permanent ID \(30006094\), title ``Expository Telescope-Conjecture Fragment,'' source code \texttt{catalogue entry}. No individual proposer is supplied; the proposed year is recorded as 2024.
\item Supplied report attribution: \emph{Arbeitsgemeinschaft: Algebraic \(K\)-Theory and the Telescope Conjecture} (2025), Oberwolfach Reports, pp. 2751--2804, DOI 10.4171/owr/2024/47. Supplied source link: \url{https://doi.org/10.4171/owr/2024/47}. The citation year 2025 and the registry's 2024 year designation are retained without reconciliation.
\item Supplied literature source: ``Oberwolfach Report 47,'' \url{https://doi.org/10.4171/OWR/2024/47}.
\item Supplied literature source: ``Hopkins and Smith, Periodicity theorem,'' \url{https://doi.org/10.2307/120991}.
\item Literature assessment checked at \texttt{2026-08-22}: this is supplied dated curation, not independently verified. The supplied status label is ``open,'' while the assessment says that the incomplete cofiber-sequence fragment admits no determinate solved or disproved status.
\item The supplied statement-verification note says that the official report identifies this extraction as background exposition and discusses the higher-height telescope conjecture as disproved. This is a supplied verification claim, not an independent source check here. The record is marked \texttt{unrecoverable}; no clean statement is supplied, and the original is retained for provenance. The record's update timestamp is \texttt{2026-09-16T12:06:53Z}.
\end{itemize}
\end{document}
